\section{Introduction}

\label{sec:introduction}

The problem of renormalization has stood at the center of modern physics
for decades.
Quantum field theory produces divergent quantities during calculation;
renormalization absorbs these into observable definitions
and yields predictions of extraordinary precision.
Yet the physical meaning of this procedure remains difficult to state.

This volume takes that difficulty as a structural clue.
We do not treat divergence merely as a mathematical artifact to be removed.
We read it as a possible boundary signal:
the trace of a descriptive window being pushed toward the limit
of the differences it can represent as finite quantities.

The name given here to this limit is the \DH{}.
It is not a wall located somewhere in spacetime,
nor a hidden object beyond physical theory.
It is a descriptive boundary that appears when a description based on
registered difference attempts to describe the absence of difference
as if it were one more observable state.

Three claims organize this volume.
First,
physical description operates on generated traces:
differences that have become stable enough to be registered,
compared,
and related.
Second,
divergence is not merely failure,
but may signal boundary contact with the limit of a given observational
window.
Third,
renormalization is not simple elimination,
but finite re-expression:
it reorganizes variables,
scales,
and reference structures so that finite relations can still be extracted.

From this perspective,
the speed-of-light limit,
the Planck-scale breakdown of spacetime continuity,
black-hole singularities,
and quantum-field divergences are not treated as unrelated failures.
Nor are they collapsed into one literal hidden object.
They are read as local or hierarchical horizon-forms:
limits that recur when a successful descriptive framework is pushed toward
the disappearance,
compression,
or overclosure of the differences on which it depends.

This framework does not propose a final unified theory.
It identifies a prior structural layer:
the condition under which registered difference,
trace formation,
and physical explanation become possible.
That condition is made explicit in the \DHP{}.

\VED{} is positioned within this framework as an open generative
description.
It does not explain existence from outside existence.
It describes why generated spacetime encounters a descriptive boundary
when it attempts to erase the differences that make its own description
possible.

The paper proceeds as follows.
Section~\ref{sec:renorm_clue} examines the unease of renormalization
and reframes it as a structural clue.
Sections~\ref{sec:past}--\ref{sec:ungenerated} establish
that physics operates on pastified traces
and that the \ungenerated{} cannot be converted into an object
of description.
Section~\ref{sec:dh} introduces the \DH{} at an intuitive level.
Sections~\ref{sec:two_modes}--\ref{sec:renorm_reinterp}
develop the reinterpretation of divergence and renormalization.
Sections~\ref{sec:replacement}--\ref{sec:unification}
re-place major theories and unsolved problems within the framework.
Section~\ref{sec:principle} states the Differential Horizon Principle.
Section~\ref{sec:ved} positions VED.
Section~\ref{sec:dynamic_barrier} gives the smallest dynamical form
of this principle as a local \DH{}.
Section~\ref{sec:conclusion} concludes.
